\section{Feature Selection Methods}
\label{sec:features}

Feature selection runs in two stages: first remove predictors that are redundant
with each other, then keep only those that carry real signal for the outcome.

\textbf{Removing collinearity by iterative VIF pruning.} From 15 numeric
candidates, instead of dropping every feature above the variance-inflation
threshold (VIF $>5$) at once, we drop the single worst offender and recompute
until every survivor clears it, because inflation is usually shared between a
pair. Several features look collinear at first (the teaching-budget ratio at
VIF $\approx 77$, individual-instruction hours at $\approx 58$, \texttt{cluster}
at $\approx 20$), but most of that collapses once the true culprit goes. Only
three drops are needed: two redundant budget ratios and the raw
\texttt{index\_value}, which duplicates the \texttt{cluster} built from it.
\texttt{cluster} itself survives once \texttt{index\_value} is gone, leaving
twelve numeric features.

\textbf{Keeping relevance by Boruta, per target.} On those 12 survivors plus
seven one-hot categorical blocks (sector, supervision, district, legal status,
education stage, locality type, year), totalling 49 encoded columns, Boruta runs
separately for each target. Using a 200-tree random forest it compares every real feature
against randomly shuffled shadow copies over up to 100 iterations, retaining one
only when its importance repeatedly exceeds the 90th percentile of the shadow
importances. A predictor survives only if its signal is dependably stronger than
noise.

The result is a \textbf{nine-feature core} confirmed for all four targets:
locality population, the nurture index, average class size, school size, and five
per-student budget ratios (tuition, perimeter, projects, purchases, transport).
Each target adds its own on top (Table~\ref{tab:boruta}). Notably
\texttt{cluster} is \emph{not} in that core: municipal status is confirmed for
three targets but rejected outright for English average grade. Five core features
are budget lines, and every target confirms at least one sector, supervision, or
district attribute, a first sign that school-level structure carries signal the
municipal cluster does not. No \texttt{year} dummy is confirmed anywhere, which
says nothing about the shape of a temporal trend but does mean no single exam
year carried enough independent signal to be retained.

Feature selection is repeated independently inside every outer training fold for
performance evaluation. Table~\ref{tab:boruta} reports the final descriptive
Boruta selection obtained by refitting on all available modelling rows after
evaluation. This final fit is used for interpretation and SHAP analysis, not for
estimating predictive performance.

\begin{table}[H]
\centering
\footnotesize
\renewcommand{\arraystretch}{1.3}
\caption{Boruta selections per target: the shared nine-feature core plus each
target's own confirmed additions. Extras are genuinely target-specific, not
cumulative. Municipal \texttt{cluster} appears as an addition for three targets
and is rejected for \texttt{english\_avg\_grade}.}
\label{tab:boruta}
\begin{tabular}{@{}>{\raggedright\arraybackslash}p{3.4cm}>{\raggedright\arraybackslash}p{8.6cm}r@{}}
\toprule
\textbf{Target} & \textbf{Additions beyond the nine-feature core} & \textbf{Total} \\
\midrule
Math grade      & \texttt{cluster}, \texttt{supervision\_Haredi} & 11 \\
Math 5-unit     & \texttt{cluster}, \texttt{district\_North}, \texttt{sector\_Jewish} & 12 \\
English 5-unit  & \texttt{cluster}, \texttt{private\_\allowbreak hours\_\allowbreak per\_student}, \texttt{sector\_Jewish}, \texttt{supervision\_Haredi} & 13 \\
English grade   & \texttt{district\_North}, \texttt{private\_\allowbreak hours\_\allowbreak per\_student}, \texttt{special\_\allowbreak ed\_share}, \texttt{sector\_Bedouin} & 13 \\
\bottomrule
\end{tabular}
\end{table}
